\chapter{System Design}
\label{ch:system-design}

% PURPOSE OF THIS CHAPTER (per ADR-055 amendment 2026-07-31):
% The thesis was designed as a three-layer system. ADR-054 froze scope to Layer 1,
% so Layers 2-3 are presented here as DESIGN ONLY, with an explicit statement that
% they were not implemented or evaluated. This keeps the architecture motivated and
% the reporting honest -- it does NOT claim unbuilt work as a contribution.
% Everything evaluated in this thesis lives in Layer 1 (Sec.~\ref{sec:design-layer1}).

\section{Design Overview}
\label{sec:design-overview}

% TODO: the three-layer decomposition and why the document-reading layer is the
% load-bearing one: every downstream layer consumes its output, so its error rate
% upper-bounds the whole system.
% Figure: three-layer architecture diagram -> thesis/figures/architecture.pdf

\section{Layer 1: Reading and Structuring}
\label{sec:design-layer1}

This is the layer implemented and evaluated in this thesis.

\subsection{Two-Stage Architecture: Reader and Structurer}
\label{sec:design-two-stage}
% TODO: separation of concerns -- a document \ac{VLM} transcribes the page image
% (the "reader"), a language model maps that transcript onto the target schema
% (the "structurer"). Contrast with the single-shot alternative (one model, image
% straight to \ac{JSON}). Both were built and compared; see Ch.~\ref{ch:results}.

\subsection{Target Schema}
\label{sec:design-schema}
% TODO: the extraction target is an \ac{EN}~16931-aligned field set (\ac{BT} codes),
% covering document header fields, party fields, the VAT breakdown group, the line-item
% group, and the payment-terms group. Rationale: the norm is the legal interoperability
% reference for German B2B invoicing, so the schema is externally defined rather than
% invented for the experiment.

\subsection{Validation and Repair}
\label{sec:design-validation}
% TODO: deterministic post-processing -- schema validation, conservative \ac{JSON}
% repair, and normalisation of representation differences (as-printed vs as-stored).
% Design principle: repair may never invent a value; an unparseable field stays absent.

\section{Layer 2 (Design Only): Compliance-Aware Knowledge Graph}
\label{sec:design-layer2}

\emph{Not implemented or evaluated within the scope of this thesis; see
Chapter~\ref{ch:limitations-future-work}.}

% TODO: the intended design -- extracted invoices become nodes/edges (supplier,
% customer, invoice, line item, tax rate), entity resolution links repeat suppliers,
% and compliance properties (mandatory-field completeness, input-tax deductibility)
% are derived per document. Why a graph: the analytically interesting questions are
% relational (per-supplier aggregation over time, duplicate detection, cross-document
% consistency), which flat retrieval expresses poorly.

\section{Layer 3 (Design Only): Analytical Query}
\label{sec:design-layer3}

\emph{Not implemented or evaluated within the scope of this thesis; see
Chapter~\ref{ch:limitations-future-work}.}

% TODO: the intended design -- natural-language questions answered over the graph,
% with a comparison of flat vector retrieval against graph-based retrieval
% (\ac{RAG} vs graph \ac{RAG}) and a per-question routing strategy.

\section{Scope of This Thesis}
\label{sec:design-scope}

% TODO (important -- the honesty anchor of the chapter):
% State plainly which layers were built and measured (Layer 1) and which remain
% design (Layers 2-3). Explain WHY: the measurement described in
% Ch.~\ref{ch:results} showed the document-reading step, not the structuring step,
% to be the binding constraint; fixing that constraint properly was judged more
% valuable than sketching two further layers on top of an unreliable foundation.
% Note that the hypotheses covering Layers 2-3 were pre-registered before this
% decision and are reported as not evaluated -- never silently dropped.
